\clearpage
\VerdierPhysicalPageAnchor{0031}
\VerdierSemanticAnchor{intro-8-etale-lefschetz}
エタール位相の枠組みにも類似の Lefschetz 公式がある [12]。
\VerdierSemanticAnchor{intro-8-frobenius-L}
この公式を、有限体上の滑らかで固有な代数曲線の Frobenius 自己準同型に
適用し、曲線にエタール位相について構成可能な層を備えることにより、一般
$L$-関数の有理性が証明された [11]。この証明では、係数層が実際に特異性
をもつ Lefschetz 公式が用いられていることに注意しよう。

\VerdierSemanticAnchor{intro-9}
\noindent\textbf{9.}\quad
位相的 Lefschetz 公式の場合、複体
\VerdierSemanticAnchor{intro-9-perfect-complex}
\[
\mathrm{RH}(X,F)
\]
が完全であることを見た。したがって
\VerdierSemanticAnchor{intro-9-rh-phi}
\[
\mathrm{RH}(\Phi)\colon
\mathrm{RH}(X,F)\longrightarrow\mathrm{RH}(X,F)
\]
の Lefschetz 数を定義できる。

\VerdierSemanticAnchor{intro-9-cohomology-objects}
この事実は、複体をそのコホモロジー対象の族で置き換えられない理由を特に
鮮明に示す。実際、自己準同型
\VerdierSemanticAnchor{intro-9-hi-phi}
\[
\mathrm{H}^i(\Phi)\colon
\mathrm{H}^i(X,F)\longrightarrow\mathrm{H}^i(X,F)
\]
のトレースは、きわめて制限的な条件のもとでしか定義できない。

\VerdierSemanticAnchor{intro-9-regular-trace}
たとえば環 $A$ の大域コホモロジー次元が有限、すなわち $A$ が正則ならば、
有限射影分解をとることで $\mathrm{H}^i(\Phi)$ のトレースを定義できる。
このとき
\VerdierSemanticAnchor{intro-9-trace-formula}
\[
\chi\bigl(\mathrm{RH}(\Phi)\bigr)
=\sum_i(-1)^i\operatorname{Tr}\bigl(\mathrm{H}^i(\Phi)\bigr)
\]
が成り立つ。

\VerdierSemanticAnchor{intro-9-nonregular-rings}
しかし、すでに見たように、とりわけエタール位相の層の理論では、
$\mathbb{Z}/p^\ell\mathbb{Z}$（$\ell\geq 2$）のような非正則環を扱わねば
ならない。

\VerdierSemanticAnchor{intro-9-k-theory}
完全複体という概念は代数幾何、すなわち K 理論と Riemann--Roch 定理
[14] においても有用である。

\VerdierSemanticAnchor{intro-10}
\noindent\textbf{10.}\quad
結論として、現在、導来圏の言葉は、これまで述べた諸理論の本質的な結果を
定式化し証明するために不可欠な道具であると思われる。

\VerdierSemanticAnchor{intro-10-additive-scope}
本論文では加法的ホモロジー代数、すなわちアーベル圏と加法的関手に限った。
Riemann--Roch 定理 [14] を扱い始めると、この枠組みでは不十分である。
\VerdierSemanticAnchor{intro-10-omissions}
また、三角圏から作られるさまざまな類群または Grothendieck 群も研究して
いない。最後に、相対ホモロジーの問題も扱わなかったが、これらも三角圏の
枠組みで処理できる [15]。
